\section{Estudio de ablación}
\label{sec:ablation}

\subsection{Diseño}

El estudio se hace sobre DistilBERT y varía dos cosas dejando todo lo demás
fijo: la arquitectura de la cabeza de clasificación y qué parte del
\textit{transformer} se entrena. Son seis configuraciones sobre los tres
datasets, es decir \num{18} corridas.

\begin{table}[h]
\centering
\small
\begin{tabular}{llll}
\toprule
\textbf{Config} & \textbf{Cabeza} & \textbf{\textit{Encoder}} & \textbf{Eje que prueba} \\
\midrule
\texttt{lineal} & $768 \to C$                & entrenable  & menos capas (ninguna oculta) \\
\texttt{h128}   & $768 \to 128 \to C$        & entrenable  & menos neuronas \\
\texttt{h768}   & $768 \to 768 \to C$        & entrenable  & referencia \\
\texttt{h2048}  & $768 \to 2048 \to C$       & entrenable  & más neuronas \\
\texttt{c2}     & $768 \to 512 \to 256 \to C$& entrenable  & más capas \\
\texttt{frz}    & $768 \to 768 \to C$        & \textbf{congelado} & congelar el \textit{transformer} \\
\bottomrule
\end{tabular}
\caption{Las seis configuraciones. $C$ es el número de clases. El diseño cubre
los \emph{dos} sentidos de cada eje --- más y menos capas, más y menos
neuronas --- porque con una sola dirección no se puede afirmar si la métrica
sube o baja de forma monótona.}
\end{table}

\paragraph{Una corrección necesaria en la configuración congelada.} La
configuración \texttt{frz} se entrena con \texttt{lr} $=10^{-3}$ en vez de
$2\times10^{-5}$. El \textit{learning rate} pequeño está pensado para ajustar
pesos ya pre-entrenados; la cabeza sobre un \textit{encoder} congelado se
aprende desde cero y con $2\times10^{-5}$ apenas avanzaría. Sin este ajuste,
\texttt{frz} perdería por el \textit{learning rate} y no por estar congelada,
que es lo que se quiere medir.

\subsection{Resultados}

\begin{figure}[h]
\centering
\includegraphics[width=0.86\textwidth]{ablation.pdf}
\caption{F1 de validación de cada configuración, por dataset. El eje llega
desde \num{0,80} a propósito: recortarlo alrededor de las diferencias haría
parecer enormes unas distancias de milésimas. Lo que se ve así es lo que
realmente ocurre --- las cinco cabezas apiladas en el mismo punto y el
\textit{encoder} congelado descolgado ---, y esa es la conclusión del
estudio.}
\label{fig:ablation}
\end{figure}

\begin{table}[h]
\centering
\small
\begin{tabular}{lcccccr}
\toprule
\textbf{Config} & \textbf{SST-2} & \textbf{AG News} & \textbf{Yelp} &
\textbf{Media} & \textbf{Params entren.} & \textbf{Tiempo (min)} \\
\midrule
\texttt{lineal} & 0,9483 & 0,9462 & 0,9603 & \best{0,9516} & 66,4 M & 15,3 \\
\texttt{h2048}  & 0,9470 & 0,9459 & 0,9601 & 0,9510 & 67,9 M & 16,3 \\
\texttt{h128}   & 0,9480 & 0,9454 & 0,9593 & 0,9509 & 66,5 M & 15,9 \\
\texttt{c2}     & 0,9476 & 0,9434 & 0,9604 & 0,9505 & 66,9 M & 15,9 \\
\texttt{h768}   & 0,9485 & 0,9427 & 0,9596 & 0,9503 & 67,0 M & 16,3 \\
\midrule
\texttt{frz}    & 0,8574 & 0,9084 & 0,9005 & 0,8887 & \best{0,6 M}  & \best{6,0} \\
\bottomrule
\end{tabular}
\caption{F1 de \textbf{validación} por configuración, ordenado por la media de
los tres datasets. El tiempo es el total de las tres corridas. La elección se
hace con esta columna; el test no participa.}
\label{tab:ablation}
\end{table}

\subsection{Lectura}

\paragraph{La cabeza de clasificación casi no importa.} Las cinco
configuraciones con el \textit{encoder} entrenable caben en \num{0,0014} de
F1. Un clasificador lineal iguala a un MLP de \num{2048} neuronas y a uno de
dos capas ocultas. Con conjuntos de validación de entre \num{6735} y
\num{20000} ejemplos, esas diferencias están dentro del ruido: ni siquiera el
orden entre ellas es estable --- \texttt{h768} es la mejor en SST-2 y la peor
en media.

Esto tiene una explicación razonable. El \textit{encoder} ya produce, en el
\textit{token} \texttt{[CLS]}, una representación de 768 dimensiones ajustada
a la tarea durante el propio \textit{fine-tuning}. Si esa representación es
linealmente separable, añadir capas por encima no aporta nada: la capacidad
que importa está en el \textit{transformer}, no en el clasificador.

\paragraph{Congelar el \textit{transformer} sí importa, y cuánto depende de la
tarea.} El \textit{encoder} congelado pierde \num{8,4} puntos de F1 en SST-2,
\num{5,8} en Yelp y solo \num{3,4} en AG News. La asimetría es informativa:
clasificar temas de noticias se resuelve casi con las características que
DistilBERT ya trae de fábrica --- es, en buena medida, detección de
vocabulario ---, mientras que detectar sentimiento en frases cortas con
negación y sarcasmo exige que el \textit{encoder} se adapte.

Vale la pena señalar el otro lado del intercambio: \texttt{frz} entrena
\num{0,6}~M de parámetros en lugar de \num{66,4}~M y tarda \num{6,0} minutos
en lugar de \num{15,3}. Para AG News, donde el costo son \num{3,4} puntos de
F1, es un intercambio que podría ser razonable si el presupuesto de
entrenamiento fuera el factor limitante.

\paragraph{Gana \texttt{lineal}, pero no por ser mejor.} Queda empatada con
las demás dentro del ruido, y es la más pequeña y la más rápida de entrenar.
Ante un empate estadístico, la configuración más simple es la elección
defendible.
